% --- LaTeX Essay Template - S. Venkatraman ---

% --- Set document class and font size ---

\documentclass[letterpaper, 11pt]{article}

% --- Package imports ---

% lipsum is just for generating placeholder text and can be removed
\usepackage{hyperref, lipsum} 

% --- Fonts (Set to Palatino or Times New Roman if desired) ---

% \usepackage{newpxtext, newpxmath}
% \usepackage{times}

% --- Page layout settiings ---

% Set page margins
\usepackage[left=1.2in, right=1.2in, top=.9in, bottom=.9in]{geometry}

% Anchor footnotes to the bottom of the page
\usepackage[bottom]{footmisc}
\usepackage{setspace}
\setstretch{1.15}
% Set line spacing
\renewcommand{\baselinestretch}{1.2}

% Use this for a one-off '*' footnote (e.g., a URL in a heading)
\newcommand{\starfootnote}[1]{%
\begingroup
\renewcommand{\thefootnote}{\fnsymbol{footnote}}%
\footnote{#1}%
\endgroup
\addtocounter{footnote}{-1}%
}

% Set spacing between paragraphs
\setlength{\parskip}{2mm}

% --- Page formatting settings ---



% --- Essay title and author ---
\title{\vspace{-1.5cm} Summary of \textit{Principal Component Analyses (PCA)-based findings in population genetic studies are highly biased and must be reevaluated}\footnote{\url{https://doi.org/10.1038/s41598-022-14395-4}}}

\date{\today}

% --- Main text ---
% --- Document starts here ---

\begin{document}
% Set link colors for labeled items and URLs (blue) 
\hypersetup{colorlinks=true, linkcolor=blue, urlcolor=blue}
\maketitle

PCA is commonly used to analyze genetic variation and infer population structure, but the authors argue that its results can be highly sensitive to data manipulation and preprocessing choices.
Through twelve test cases and human population datasets, they demonstrate that PCA plots can be manipulated to produce different, even contradictory, conclusions.

Also, PCA is commonly used downstream as well as other fields. If the principal components are artifacts of data manipulation,
then including them for analyses will introduce bias.

The paper calls for caution in interpreting PCA results and recommends that researchers test the robustness of their findings to different preprocessing choices.

\end{document}